\begin{statementbox}
  \textbf{\textcolor{CStatement}{条件}}
  \begin{description}[style=nextline, leftmargin=2.8em, labelsep=0.5em, itemsep=0.35em]
    \item[\textbf{\textcolor{CStatement}{条件 1（标准方程与点在曲线上）：}}]
      抛物线方程为 $y^2=2px$（$p>0$），且点 $P(x_0,y_0)$ 在该抛物线上，因而 $x_0\ge 0$。
  \end{description}

  \textbf{\textcolor{CStatement}{结论}}
  \begin{description}[style=nextline, leftmargin=2.8em, labelsep=0.5em, itemsep=0.45em]
    \item[\textbf{\textcolor{CStatement}{结论一（向右开口的焦半径公式）：}}]
      \textbf{\textcolor{CStatement}{直观描述：}}
      焦半径长度等于点的横坐标与焦点到顶点距离之和（后者是 $p/2$）。
      \[
        \boxed{PF=x_0+\frac p2}.
      \]
      焦点 $F(\tfrac p2,0)$，准线 $l:x=-\tfrac p2$。
    \item[\textbf{\textcolor{CStatement}{常见变形（向上开口）：}}]
      若 $x^2=2py$（$p>0$），$P(x_0,y_0)$ 在曲线上，则
      \[
        \boxed{PF=y_0+\frac p2}.
      \]
  \end{description}

  \medskip
  \textbf{\textcolor{CStatement}{统一记法（四个方向只记一条规则）}}

  \textbf{沿开口方向，从顶点到 $P$ 的投影距离 $t$，再加焦点到顶点的距离 $p/2$：}
  \[
    \boxed{PF=t+\frac p2},\qquad t\ge 0.
  \]
  对于以下四种标准方程，$P$ 均须在相应曲线上且 $p>0$：
  \[
    \begin{array}{c|c|c}
      \text{方程}&t&\text{焦半径}\\ \hline
      y^2= 2px& x_0  & PF=x_0+\frac p2\\
      y^2=-2px&-x_0  & PF=-x_0+\frac p2\\
      x^2= 2py& y_0  & PF=y_0+\frac p2\\
      x^2=-2py&-y_0  & PF=-y_0+\frac p2
    \end{array}
  \]
  \textbf{一句话记忆：}看开口，取对应坐标的非负长度，再加 $p/2$。

  \medskip
  \textbf{\textcolor{CStatement}{参数提醒}}
  $y^2=2px$ 中 $p$ 是焦点到准线的距离；焦点到顶点为 $p/2$；通径长为 $2p$，半通径长为 $p$。
  若课本写 $y^2=4ax$（$a>0$），则 $a=p/2$，所以 $PF=x_0+a$。
\end{statementbox}
